\chapter*{Resumen}

El seguimiento multiobjeto (\emph{multi-object tracking}, MOT) de jugadores
en vídeo de fútbol es un paso previo necesario para cualquier análisis
automático del juego como mapas de calor, distancia recorrida, posesión de
balón, pero elegir un algoritmo concreto para esa tarea exige compararlo
sobre datos reales de fútbol y no solo sobre los benchmarks genéricos de
seguimiento de personas para los que se diseñó originalmente. Este Trabajo
de Fin de Grado aborda esa comparación sobre SoccerNet-Tracking, un
\emph{dataset} público de vídeo de fútbol con detecciones precomputadas, y
propone un sistema completo que va desde el seguimiento hasta la
extracción de metadatos deportivos interpretables a partir de las
trayectorias resultantes.

El sistema compara dos algoritmos de
asociación clásicos, ByteTrack y OC-SORT, evaluados con las métricas
estándar de TrackEval (HOTA, MOTA, IDF1) sobre los conjuntos de
entrenamiento y de test. Sobre la configuración óptima de cada algoritmo,
obtenida mediante un estudio sistemático de hiperparámetros, se construye
un segundo pipeline de extracción de metadatos: clasificación de equipos
por color de camiseta mediante K-means, mapas de calor de posiciones,
distancia y velocidad reales a partir de una homografía calibrada
manualmente, y posesión de balón por proximidad. Los resultados, junto con
vídeos anotados, se exponen finalmente en una aplicación web.

El conjunto de datos procesado reúne 106 secuencias de vídeo (57 de
entrenamiento, 49 de test) de 700 fotogramas cada una y con
detecciones precomputadas de jugadores, árbitros y balón. Con la métrica
HOTA(0) (evaluada en un único umbral),
ByteTrack alcanza 81,43 en test frente a 80,92 de OC-SORT; sin embargo,
con la métrica HOTA integrada estándar el resultado se invierte con
claridad en ambos conjuntos (80,23 frente a 73,27 en entrenamiento; 75,43
frente a 70,73 en test), porque OC-SORT reutiliza las detecciones del
\emph{dataset} sin modificarlas mientras que ByteTrack sí las interpola en
su segunda ronda de asociación. Esta discrepancia, documentada y explicada
en detalle en la memoria, es uno de los resultados más relevantes del
trabajo: dos formas igual de legítimas de medir el mismo sistema pueden
llevar a conclusiones opuestas sobre qué algoritmo es mejor.

\vspace{.5cm}

\textbf{Palabras clave:} SoccerNet, Seguimiento Multiobjeto, Visión
Artificial, ByteTrack, OC-SORT, HOTA, Análisis Deportivo.
